\begin{thema}{Γ}
Θεωρούμε την εξίσωση $e^x = a x$ με παράμετρο $a \in \mathbb{R}$.
\begin{erwthma}
  \item Να δείξετε ότι η συνάρτηση $g(x) = \dfrac{e^x}{x}$ με $x>0$ παρουσιάζει ελάχιστο στο $x=1$ και να βρείτε την ελάχιστη τιμή της.
  \item Αποδείξτε ότι, για $a > e$, η ευθεία $y = a x$ τέμνει την καμπύλη $y = e^x$ σε ακριβώς δύο σημεία στο πρώτο τεταρτημόριο. 
  \item Για $a \le e$, πόσες λύσεις στο $(0, +\infty)$ μπορεί να έχει η εξίσωση $e^x = a x$; Να αιτιολογήσετε την απάντησή σας.
  \item Να υπολογίσετε το όριο $\lim\limits_{x \to 1} \dfrac{g(x) - e}{(x-1)^2}$.
  \item Να αποδείξετε ότι η $g$ είναι κυρτή στο $(0, +\infty)$ και ότι για $1 < \alpha < \beta$ ισχύει $g'(\alpha) < \dfrac{g(\beta) - g(\alpha)}{\beta - \alpha} < g'(\beta)$.
\end{erwthma}
\end{thema}
